\originalpreface
本册以会初等微积分、尚未系统学习量子力学的自学者为读者. 复数振幅、概率密度、波函数和算符在首次使用时解释. 单变量微积分和常微分方程的基本运算作为先修; 希尔伯特空间的严格谱理论没有假定已经掌握. 无界算符的论证在文中说明所用边界和共同定义域, 不把形式计算当作完整的泛函分析证明.

材料主轴为 Barton Zwiebach 的 8.04 Quantum Physics I (Spring 2016)、8.05 Quantum Physics II (Fall 2013) 和 8.06 Quantum Physics III (Spring 2018). 本册是中文重新组织和补充推导的非官方改编, 没有 MIT 官方中文或认可的含义. 原始文件保留在仓库量子力学来源目录. 原件和中文改编遵循 \href{https://creativecommons.org/licenses/by-nc-sa/4.0/}{CC BY-NC-SA 4.0}, 署名 MIT OpenCourseWare、Barton Zwiebach 及原件另署的贡献者. 本册新增解释、例题和答案标为独立整理. 不使用 MIT 校徽或宣传性标志.

本册现有19章理论, 39份作业、测试与期末文档中227道可用大题的全部小问. 原17理论章与19考核篇保留, 新增谱与变分、含时辐射两章, 第8章接入线代补充, 20份新增作业共115题. 两张公式表另归档. 题解均为\textbf{AI独立解答}, 所选OCW页面没有取得官方答案. 23份先前被整份暂缓的原件现全部恢复; 本次按完整MIT任务逐题补入, 商业教材只有题号而没有题面的10道不编造. 改编信用保留, 表格与数值图自主推导和绘制. 独立模型复核不称为教师答案, 本册也不是三课每页原件的全文翻译.

\noindent\textbf{8.04 课程安排.}
每周两次90分钟讲授和两次60分钟习题课; 官方先修是8.03振动与波, 成绩至少C. 三部分为量子基本概念、一维势中的运动、一维散射与中心势. 两次课内考试和三小时期末. 总评为考试平均75\%、作业25\%; 考试平均中两次测试各25\%、期末50\%, 所以折合总评分别18.75\%、18.75\%、37.5\%. 作业最低或缺交的一套不计入平均. 2016页面实际提供2015秋测试, 年份不能合并.

\noindent\textbf{8.05 课程安排.}
每周两次90分钟讲授和两次60分钟习题课; 先修8.04, 至少C. 大纲包括量子力学一般结构、动力学、二态系统、角动量与自旋、径向方程、角动量合成、全同粒子. 一次90分钟闭卷期中和三小时期末. 总评: 期中25\%、期末45\%、作业30\%. 官网发布11套作业, 以及2011--2013年的期中与期末试卷和两张公式表.

\noindent\textbf{8.06 课程安排.}
每周两次90分钟讲授和两次60分钟习题课; 先修8.05, 至少C. 大纲以定态微扰、氢精细结构、WKB、含时微扰、带电粒子、绝热近似、散射和全同粒子展开. 总评: 作业25\%、期中15\%、学期论文25\%、期末35\%. 论文成绩须至少C才能通过, 另有提案、初稿与同伴编辑环节. 课纲述九套计分作业和末周一套不计分作业, OCW作业页只有九套. 90分钟期中和三小时期末\textbf{明确不向OCW公开}; 第5章带电粒子讲义也未提供. 本册没有伪造这些题面或章节原件.

下列课历保留官方讲次与主题顺序; 主题用中文表述, 完整官网链接见卷末来源对应. 作业序号和考试位置保留; 官网未给具体日历日期, 不另造日期.

\noindent\textbf{8.04 课历.}
\begin{longtable}{p{.10\textwidth}p{.67\textwidth}p{.14\textwidth}}
\toprule 讲次 & 主题 & 考核节点\\\midrule\endhead
1--2 & 线性、复数、随机性、叠加与纠缠; 双缝、干涉仪、炸弹探测 & \\
3--5 & 光电效应、康普顿散射、物质波; 伽利略变换、群速度; 动量与位置算符、薛定谔方程 & 作业1、2\\
6--9 & 归一化、概率流、厄米算符; 波包与傅里叶变换; 期望值、Ehrenfest方程; 测量和不确定性 & 作业3、4\\
10--12 & 定态、圆周粒子、无限和有限势阱; 转折点、半经典近似 & 测试1; 作业5\\
13--15 & 打靶法; $\delta$ 势、节点; 谐振子微分方程和升降算符 & 作业6\\
16--18 & 散射、反射和透射; 共振透射; 相移、Levinson定理 & 作业7; 测试2\\
19--21 & 共振、复波数; 角动量代数、勒让德函数; 径向方程与氢原子 & 作业8、9\\
22--24 & 氢原子尺度、级数与量子化; 维里关系、轨道与转折点; 自旋 & 作业10; 期末\\
\bottomrule
\end{longtable}
\noindent\textbf{8.05 课历.}
\begin{longtable}{p{.12\textwidth}p{.65\textwidth}p{.14\textwidth}}
\toprule 讲次 & 主题 & 作业\\\midrule\endhead
1--3 & 波动力学 & 1\\
3--4 & 自旋$1/2$、bra、ket与算符 & 2\\
5--8 & 向量空间与算符 & 3、4\\
9 & Dirac记号 & 5\\
10--11 & 不确定性与相容观测量 & 6\\
12--16 & 量子动力学 & 7\\
16--18 & 二态系统 & 8\\
18--20 & 多粒子态与张量积 & 9\\
20--23 & 角动量 & 10\\
24--26 & 角动量合成 & 11\\
\bottomrule
\end{longtable}
\noindent\textbf{8.06 课历.}
\begin{longtable}{p{.12\textwidth}p{.65\textwidth}p{.14\textwidth}}
\toprule 讲次 & 主题 & 节点\\\midrule\endhead
L1--L3 & 定态与简并微扰 & 作业1\\
L4--L6 & 氢精细结构、Zeeman效应、半经典近似引入 & 作业2\\
L7--L9 & WKB、Airy函数和连接式; 含时微扰 & 作业3、4\\
L10--L11 & Fermi黄金规则和谐波跃迁 & 作业5; E1\\
L12--L14 & 氢电离、光与原子、带电粒子与电磁场 & 论文提案\\
L15--L18 & 绝热近似、Berry相位、分子 & 作业6、7; 论文初稿\\
L19--L22 & 散射、全同粒子引入 & 作业8\\
L23--L24 & 全同粒子、论文返还与期末复习 & E2\\
\bottomrule
\end{longtable}
